% ЗАКЛЮЧЕНИЕ — без номер. Ориентировъчен обем: 2–3 стр.
% По А.3 съдържа: обобщение на характеристиките и особеностите на решението,
% предимствата И НЕДОСТАТЪЦИТЕ му, резултатите от експериментите с тяхната
% инженерна трактовка, и предложения за по-нататъшна работа.
\section*{ЗАКЛЮЧЕНИЕ}
\label{sec:conclusion}
\addcontentsline{toc}{section}{ЗАКЛЮЧЕНИЕ}

В хода на дипломното проектиране е изградена работеща система, която прави
съдържанието на предварителните товарни таблици на Walltopia достъпно през браузър,
запазва изчисленията като проекти на конкретен потребител и позволява от тях да бъде
зададено запитване към инженерния отдел. По задачите от Увода, в същия ред: описана е
проблемната област, изведени са петте недостатъка на съществуващото решение и от тях
единадесет функционални и осем нефункционални изисквания (Раздел~\ref{sec:domain});
проектирани са архитектурата на системата и базата данни, като изборът между
разгледаните варианти е обоснован срещу предварително записани критерии
(Раздел~\ref{sec:architecture}); архитектурата е реализирана като услуга на Node.js и
уеб клиент без стъпка на компилация, а реализацията е проверена срещу три различни по
вид еталона с представени и анализирани резултати (Раздел~\ref{sec:implementation});
и е изготвено ръководство на потребителя, което описва работата със системата от
гледна точка на човека, който я използва (Раздел~\ref{sec:userguide}).

\textbf{Характеристики и особености на решението.} Системата има три свойства, които
я отличават от обичайното уеб приложение със същото предназначение. Първо,
изчислението се извършва изцяло в клиента върху вграден набор от данни, поради което
калкулаторът остава работоспособен и когато базата от данни е недостъпна; това не е
следствие от реализацията, а е заложено в реда на монтиране на маршрутите. Второ,
договорът на приложния програмен интерфейс приема маркера за удостоверяване както от
бисквитка, така и от заглавие, което го държи отворен за всеки консуматор извън
браузъра, без уеб клиентът да прави компромис с начина, по който пази собствената си
сесия. Трето, коректността на изчислителното ядро е измерена срещу независим еталон,
а не приета на доверие.

\textbf{Предимства.}
\begin{enumerate}[topsep=2pt,itemsep=2pt,leftmargin=*]
\item \textbf{Проверимост.} Съществуването на еталон превръща основното изискване в
      измерима величина. Резултатът от §\ref{subsec:oracle} показва, че това не е
      формалност: правдоподобната алтернатива на приетото решение би нарушила
      изискването при около 47 \% от входните комбинации, и то незабележимо.
\item \textbf{Работоспособност при частичен отказ.} Отпадането на базата от данни
      отнема профилите и проектите, но не и калкулатора, който е основната функция.
\item \textbf{Ниска цена на доставянето.} Уеб клиентът се публикува чрез копиране на
      директория, без инструментална верига и без стъпка на компилация, и се обслужва
      от същата услуга, която носи приложния програмен интерфейс.
\item \textbf{Проследимост.} Всяко изчисление може да бъде запазено с име, етикети и
      произволни свойства, а запитването към инженерния отдел носи изчислителния
      контекст автоматично.
\end{enumerate}

\textbf{Недостатъци и ограничения.} Изчислителният слой на уеб клиента не е отделен
файл и не може да бъде изпитан изолирано (§\ref{subsec:limitations}); това е пряката
цена на отказа от стъпка на компилация и е най-същественият структурен компромис в
работата. Наборът с данните съществува в две копия, чието съвпадение се установява с
проверка, а не се гарантира по конструкция, а обобщението на проекта се съхранява
отделно от входа и остарява при промяна в товарните таблици. Проверката на интерфейса
измерва едно свойство, не качество: половината от установените класове дефекти
(§\ref{subsec:taxonomy}) са открити чрез оглед, а не от нея. Тя е проведена върху един
двигател за изобразяване и една версия на еталонната работна книга, като нито едното,
нито другото се пренася автоматично, а изпитанието на приложния програмен интерфейс
използва заместител на базата от данни. Накрая, документът с описанието на приложния
програмен интерфейс е остарял в две точки спрямо реализацията.

\textbf{Инженерна трактовка на резултатите.} Двата основни експеримента водят до един
и същ извод, макар да са от различни области. При числата (§\ref{subsec:oracle}) и при
разположението (§\ref{subsec:viewportres}) най-опасният дефект не е този, който проваля
системата, а този, който я оставя да изглежда изправна: имперска стойност, различна с
0{,}1 lb, и съдържание, което липсва, без да е съобщено, че липсва. И в двата случая
откриването е било възможно само чрез механично измерване на цялото пространство, а
не чрез проверка на подбрани случаи. Това е основанието проверката да бъде описана
подробно (§\ref{subsec:method}) преди резултатите от нея.

\textbf{Насоки за по-нататъшна работа.}
\begin{enumerate}[topsep=2pt,itemsep=1pt,leftmargin=*]
\item \textbf{Един източник за товарните таблици}: двете копия да се произвеждат от
      един файл при доставянето, с което ограничението отпада по конструкция.
\item \textbf{Отделяне на изчислителния слой на уеб клиента} чрез модули на браузъра;
      те се поддържат без подготовка, тоест без стъпка на компилация.
\item \textbf{Включване на проверката по §\ref{subsec:howcheck} в непрекъсната
      интеграция.} Поведението от §\ref{subsec:taxonomy}, при което една поправка
      разкрива следващия дефект, прави еднократното изпълнение недостатъчно.
\item \textbf{Разширяване на проверката към втори двигател за изобразяване}, за да се
      провери дали систематизацията от §\ref{subsec:taxonomy} е свойство на задачата
      или на конкретния браузър.
\item \textbf{Извличане на данните от работната книга като част от процеса}, така че
      нова редакция да поражда ново изпълнение на сравнението от §\ref{subsec:oracle}.
\item \textbf{Изнасяне на историята на запитванията в интерфейса}: записите се създават,
      но входната точка за четенето им не е реализирана.
\item \textbf{Проверка за достъпност} по WCAG 2.2~\cite{wcag22}, която използва същия
      стенд и не е обхваната в тази работа.
\item \textbf{Синхронизиране на мобилното приложение с уеб реализацията.} В
      хранилището съществува отделно мобилно приложение (\code{mobile/}), извън
      обхвата на тази работа (Увод, §\ref{sec:introduction}), което към момента
      изостава от описаната тук уеб реализация, включително от конфигуратора на
      закрепването (§\ref{subsec:attachimpl}). Довеждането му до същото функционално
      ниво е самостоятелна задача, а не разширение на настоящата работа, защото изисква
      собствена проверка на съответствието с еталона за платформата на мобилното
      приложение.
\end{enumerate}
